\section{Przegląd dostępnych narzędzi i usług Speech-to-Text} Współczesny rynek technologii automatycznego rozpoznawania mowy (ASR / STT) oferuje szeroki wachlarz rozwiązań, które można zaklasyfikować do trzech głównych kategorii architektonicznych: komercyjnych usług chmurowych (ang. \textit{Cloud API}), otwartych i lokalnych modeli głębokich (ang. \textit{Open-Source / On-Device}) oraz natywnych mechanizmów przeglądarkowych. Wybór odpowiedniego narzędzia w projektach programistycznych wymaga przeanalizowania kompromisu pomiędzy dokładnością transkrypcji, czasem odpowiedzi systemu (latencją), kosztami eksploatacji oraz wymogami dotyczącymi prywatności danych. Do grupy wiodących dostawców komercyjnych API chmurowych należą platformy Google Cloud Speech-to-Text, Microsoft Azure Speech Service oraz Amazon Transcribe. Usługi te działają w modelu serwerowym (\textit{Software as a Service} -- SaaS), w którym akwizycja i wstępne przetwarzanie audio zachodzą po stronie klienta, natomiast właściwe dekodowanie akustyczne i językowe realizowane jest na infrastrukturze dostawcy chmury. Rozwiązania te oferują wysokie wsparcie dla przetwarzania strumieniowego w czasie rzeczywistym oraz stałe aktualizacje modeli bez konieczności angażowania zasobów sprzętowych użytkownika końcowego. Alternatywą dla usług chmurowych są samodzielne modele uczenia głębokiego uruchamiane lokalnie (\textit{On-Device} / \textit{Self-Hosted}), wśród których standardem stał się model OpenAI Whisper oraz jego zoptymalizowane warianty (np. WhisperX) \cite[s.~2--3]{ferraro2023benchmarking}. Modele te pozwalają na pełną autonomię przetwarzania, brak opłat subskrypcyjnych za pojedyncze zapytania oraz gwarancję prywatności danych audio, wymagają jednak odpowiednich zasobów sprzętowych (pamięci VRAM oraz wydajnych akceleratorów GPU) po stronie serwera lub urządzenia brzegowego \cite[s.~18]{pudo2024optimizing}. Ostatnią kategorię stanowią interfejsy natywne, reprezentowane przez standard W3C Web Speech API, dostarczający wbudowane w przeglądarki internetowe API do rozpoznawania mowy (\texttt{SpeechRecognition}) \cite[s.~1]{mdnWebSpeechAPI}. Choć rozwiązanie to wyeliminowuje potrzebę zewnętrznych integracji i opłat, jego działanie jest silnie uzależnione od silnika dostarczanego przez producenta danej przeglądarki i systemu operacyjnego.

\subsection{Kryteria porównawcze narzędzi STT} Ocena przydatności i efektywności poszczególnych rozwiązań STT w aplikacjach transkrypcyjnych wymaga zdefiniowania ujednoliconych kryteriów ilościowych oraz jakościowych. W niniejszej pracy przyjęto zestaw następujących wskaźników porównawczych:
    \begin{itemize}
        \item \textbf{Dokładność transkrypcji (WER oraz CER):} Podstawową miarą jakości rozpoznawania mowy jest wyrazowa stopa błędów (ang. \textit{Word Error Rate} -- WER), definiowana jako stosunek sumy podstawień ($S$), usunięć ($D$) i wstawień ($I$) do całkowitej liczby słów w tekście wzorcowym ($N$) \cite[s.~330]{columbia_asr}:
        \begin{equation}
            \text{WER} = \frac{S + D + I}{N} \cdot 100\%
        \end{equation}
        Uzupełnieniem WER jest znakowa stopa błędów (ang. \textit{Character Error Rate} -- CER), obliczana analogicznie na poziomie poszczególnych liter i znaków, co pozwala na precyzyjniejszą ocenę błędów w językach bogatych fleksyjnie.
        \item \textbf{Opóźnienia i wydajność (Real-Time Factor -- RTF):} Czas przetwarzania sygnału audio opisywany jest za pomocą współczynnika czasu rzeczywistego (ang. \textit{Real-Time Factor} -- RTF), stanowiącego stosunek czasu trwania procesu transkrypcji ($T_{\text{proc}}$) do długości nagrania audio ($T_{\text{audio}}$):
        \begin{equation}
            \text{RTF} = \frac{T_{\text{proc}}}{T_{\text{audio}}}
        \end{equation}
    Wartość $\text{RTF} < 1$ oznacza, że system przetwarza nagranie szybciej niż trwa jego odtwarzanie w czasie rzeczywistym.
    \item \textbf{Koszty użytkowania i modele rozliczeń:} W przypadku API chmurowych analizowane są darmowe limity miesięczne (ang. \textit{Free Tier}) oraz stawki cenowe naliczane za każdą rozpoczętą minutę przetworzonego nagrania audio.
    \item \textbf{Autonomia i praca offline:} Ocena możliwości uruchomienia silnika rozpoznawania mowy bez konieczności utrzymywania aktywnego połączenia internetowego i przesyłania strumienia audio do zewnętrznych serwerów.
    \item \textbf{Dostępność i specyfika języka polskiego:} Odporność modeli na charakterystyczne cechy języka polskiego, takie jak bogata odmiana końcówek (fleksja), swobodny szyk zdań, zapis polskich znaków diakrytycznych oraz spójne formatowanie danych liczbowych i godzinowych. \end{itemize}


\subsection{Google Speech-to-Text API} Google Cloud Speech-to-Text (GCP STT) to usługa transkrypcyjna udostępniana w chmurowym modelu SaaS (\textit{Software as a Service}), wykorzystująca głębokie sieci neuronowe opracowane przez firmę Google. Platforma zapewnia obsługę ponad 125 języków i wariantów regionalnych (w tym języka polskiego) oraz umożliwia przetwarzanie mowy w trybie wsadowym oraz strumieniowym w czasie rzeczywistym z wykorzystaniem protokołów REST oraz gRPC \cite{google_stt_v2_doc}. Dostawca udostępnia dojrzałą wersję interfejsu API v1 oraz nowszą generację v2 (\texttt{Google.Cloud.Speech.V2}), która wprowadza rozbudowane klasy konfiguracyjne (np. \texttt{RecognitionFeatures}) i nowe modele architektury (m.in. rodzinę modeli \textit{Chirp}) \cite{google_stt_v2_doc}. Wśród kluczowych funkcji programistycznych usługa oferuje automatyczną interpunkcję (\textit{automatic punctuation}), rozpoznawanie dyktowanych symboli i emotikonów, diarizację mówców (\textit{speaker diarization}) oraz adaptację słownikową (\textit{speech adaptation} / \textit{PhraseSet}), pozwalającą na dostarczanie kontekstowych list słów kluczowych i terminologii branżowej \cite{google_stt_v2_doc}. Model finansowy usługi w oficjalnej dokumentacji Google Cloud różnicuje stawki w zależności od wersji API, trybu przetwarzania oraz zgody na prywatność danych \cite{google_stt_pricing}: \begin{itemize} \item \textbf{Wersja API v1:} Oferuje darmowy pakiet początkowy (ang. \textit{Free Tier}) obejmujący 60 minut transkrypcji miesięcznie. Po jego przekroczeniu podstawowa stawka wynosi $0{,}016$\,\$ za minutę przy wyrażeniu zgody na rejestrowanie nagrań w celu doskonalenia modeli Google (ang. \textit{with data logging}) lub $0{,}024$\,\$/min wariancie zapewniającym pełną poufność danych (ang. \textit{without data logging}) \cite{google_stt_pricing}. \item \textbf{Wersja API v2:} Wprowadza elastyczny model opłat zależny od miesięcznego wolumenu (rozpoczynający się od $0{,}016$\,\$/min dla standardowych żądań do $500\,000$ minut) oraz zoptymalizowaną opcję przetwarzania wsadowego (ang. \textit{Dynamic Batch Recognition}), w której stawka zostaje zredukowana do $0{,}003$\,\$ za minutę audio \cite{google_stt_pricing}. \end{itemize} Architektura Google STT charakteryzuje się bardzo wysoką skalowalnością i niskim współczynnikiem opóźnień (RTF), jednak ze względu na realizację dekodowania po stronie infrastruktury chmurowej wymusza stałą i stabilną łączność z siecią Internet.

\subsection{Microsoft Azure Speech Service} Microsoft Azure Speech Service (wchodzący w skład pakietu Azure AI Services) to komercyjna usługa chmurowa wykorzystująca zaawansowane modele głębokich sieci neuronowych do transkrypcji mowy \cite{azure_speech_overview}. Integracja z usługą realizowana jest poprzez dedykowany pakiet programistyczny \textit{Speech SDK}, interfejs wiersza poleceń \textit{Speech CLI} lub REST API \cite{azure_speech_overview}. Platforma udostępnia trzy podstawowe tryby przetwarzania audio:
\begin{itemize}
    \item \textbf{Transkrypcja w czasie rzeczywistym (ang. \textit{Real-time transcription}):} Natychmiastowe rozpoznawanie mowy z mikrofonu lub strumienia audio z błyskawicznym zwracaniem wyników pośrednich.
    \item \textbf{Szybka transkrypcja (ang. \textit{Fast transcription}):} Synchroniczne przetwarzanie plików audio ze stałą latencją dla pojedynczych nagrań z wykorzystaniem zoptymalizowanych punktów końcowych.
    \item \textbf{Transkrypcja wsadowa (ang. \textit{Batch transcription}):} Asynchroniczna analiza dużych wolumenów nagrań zgromadzonych w magazynach danych.
\end{itemize} Usługa dostarcza bogaty zestaw funkcji wspomagających jakość transkrypcji, takich jak automatyczna interpunkcja, diarizacja pozwalająca na wyodrębnienie i rozróżnienie wielu mówców, rozpoznawanie dwukanałowe (ang. \textit{multichannel transcription}) dla nagrań stereo oraz automatyczna detekcja języka \cite{azure_speech_overview}. Ponadto moduł \textit{Custom Speech} oraz listy fraz (\textit{phrase lists}) umożliwiają dostosowanie modeli bazowych do unikalnego słownictwa branżowego i specyficznych warunków akustycznych \cite{azure_speech_overview}. Model rozliczeniowy opiera się na rzeczywistym czasie przetworzonego audio naliczanym z dokładnością do 1 sekundy \cite{azure_speech_pricing}. W ramach darmowej warstwy (\textit{Free F0}) Microsoft oferuje 5 darmowych godzin audio miesięcznie dla transkrypcji w czasie rzeczywistym \cite{azure_speech_pricing}. W modelu \textit{Pay-As-You-Go} standardowa transkrypcja w czasie rzeczywistym kosztuje $1{,}00$\,\$ za godzinę audio (ok. $0{,}0167$\,\$/min), transkrypcja wsadowa wynosi $0{,}18$\,\$ za godzinę, natomiast funkcje rozszerzone (takie jak ciągła detekcja języka czy diarizacja) rozliczane są jako dodatek w cenie $0{,}30$\,\$ za godzinę za każdą funkcję \cite{azure_speech_pricing}.

\subsection{Amazon Transcribe} Amazon Transcribe to w pełni zarządzana usługa automatycznego rozpoznawania mowy (ASR) udostępniana w chmurowej infrastrukturze Amazon Web Services (AWS) \cite{amazon_transcribe_features}. Platforma umożliwia przekształcanie sygnału mowy na tekst zarówno w trybie wsadowym (asynchroniczna analiza gotowych plików audio), jak i w trybie strumieniowym w czasie rzeczywistym z wykorzystaniem bezpiecznych połączeń sieciowych \cite{amazon_transcribe_features}. System wspiera szeroki wachlarz formatów i kontenerów audio/wideo (m.in. WAV, MP3, FLAC, MP4) i jest bezpośrednio zintegrowany z magazynem danych Amazon S3 oraz systemem monitorowania AWS CloudWatch \cite{amazon_transcribe_features}. Wśród kluczowych funkcji programistycznych wspierających jakość i analizę transkrypcji znajdują się \cite{amazon_transcribe_features}: \begin{itemize} \item \textbf{Automatyczna interpunkcja i normalizacja cyfr:} Automatyczne formatowanie tekstu i konwersja wypowiadanych liczb na ich postać cyfrową. \item \textbf{Generowanie znaczników czasu i diarizacja:} Precyzyjne wyznaczanie znaczników czasu (ang. \textit{timestamps}) na poziomie pojedynczych słów oraz automatyczna diarizacja i separacja kanałów audio dla nagrań wielokanałowych/stereo. \item \textbf{Automatyczna identyfikacja języka:} Samoczynne wykrywanie języka dominującego w nagraniu lub obsługa płynnego przełączania się między wieloma językami bez konieczności ręcznego definiowania kodu językowego. \item \textbf{Dostosowanie i adaptacja dziedzinowa:} Tworzenie własnych słowników (\textit{Custom Vocabulary}) dla specyficznej terminologii branżowej oraz trenowanie dedykowanych modeli językowych (\textit{Custom Language Models} -- CLM) na własnych korpusach tekstowych. \item \textbf{Bezpieczeństwo i anonimizacja danych:} Filtrowanie niepożądanego słownictwa (\textit{vocabulary filtering}), detekcja toksyczności wypowiedzi oraz automatyczna redakcja i anonimizacja wrażliwych danych osobowych (\textit{PII redaction}). Szyfrowanie danych w spoczynku (Amazon S3 / AWS KMS) oraz w transmisji (TLS 1.2). \end{itemize} Usługa udostępnia również dedykowane interfejsy API dostosowane do specyficznych domen branżowych, takie jak \textit{Amazon Transcribe Call Analytics} (rozszerzony o generatywne podsumowania rozmów i analizę sentymentu) oraz \textit{Amazon Transcribe Medical} (przystosowany do dyktowania i konwersacji medycznych) \cite{amazon_transcribe_features}. Model rozliczeniowy Amazon Transcribe opiera się na rzeczywistym czasie przetworzonego audio naliczanym sekundowo bez minimalnego czasu żądania \cite{amazon_transcribe_pricing}. W ramach bezpłatnej warstwy \textit{AWS Free Tier} dostawca oferuje pakiet 60 darmowych minut miesięcznie przez pierwsze 12 miesięcy od utworzenia konta \cite{amazon_transcribe_pricing}. W standardowym modelu opłat opłaty rozliczane są według stawek zależnych od wolumenu (np. transkrypcja wsadowa od $0{,}006$\,\$/min, transkrypcja strumieniowa od $0{,}01$\,\$/min dla regionu US East), natomiast funkcje rozszerzone (np. anonimizacja PII, modele CLM czy moduły analityczne) stanowią dodatkowo płatne komponenty \cite{amazon_transcribe_pricing}.

% dotąd sprawdzone
\subsection{OpenAI Whisper (on-device i w chmurze)} OpenAI Whisper to wielozadaniowy, otwarty system automatycznego rozpoznawania mowy oparty na architekturze sekwencja-do-sekwencji (Transformer), wytrenowany w reżimie słabego nadzoru na ogromnym zbiorze 680\,000 godzin wielojęzycznych i wielozadaniowych danych audio zgromadzonych z sieci Internet \cite[s.~1, 3--4]{radford2023robust}. Architektura ta przetwarza sygnał audio w 30-sekundowych okienkach przeliczanych na log-spektrogramy Mel, łącząc rozpoznawanie mowy z detekcją języka, automatyczną interpunkcją oraz możliwością bezpośredniej translacji na język angielski \cite[s.~3--5]{radford2023robust}. Model oferowany jest w dwóch zasadniczych paradygmatach wdrożeniowych: lokalnym (\textit{On-Device / Self-Hosted}) oraz chmurowym (\textit{Hosted API}) \cite{openai_file_transcription}, \cite[s.~18]{pudo2024optimizing}. \begin{itemize} \item \textbf{Wariant lokalny (On-Device / Open-Source):} OpenAI udostępnia nieodpłatnie kod źródłowy oraz wagi modeli w kilku wariantach rozmiarowych — od wersji \texttt{tiny} (39M parametrów) i \texttt{base} (74M), przez \texttt{small} (244M) i \texttt{medium} (769M), aż po najnowszą generację \texttt{large-v3} (1{,}55B parametrów) \cite[s.~5]{radford2023robust}. Uruchomienie modelu w środowisku lokalnym gwarantuje $100\%$ prywatności danych oraz zero opłat subskrypcyjnych za przetwarzane minuty nagrań, wymaga jednak odpowiednich zasobów sprzętowych (pamięci VRAM od 1\,GB do ponad 10\,GB oraz akceleracji GPU) \cite[s.~18]{pudo2024optimizing}. W środowisku programistycznym rozwijane są dodatkowo zoptymalizowane silniki, takie jak \textbf{WhisperX}, które integrują wstępną detekcję aktywności głosu (VAD), wymuszone wyrównanie fonetyczne (\textit{forced alignment}) oraz przetwarzanie wsadowe (\textit{batching}), co drastically skraca czas przetwarzania (RTF) i eliminuje problem pętli powtórzeń na długich nagraniach \cite[s.~2]{ferraro2023benchmarking}. \item \textbf{Wariant chmurowy (OpenAI API):} Dostępny poprzez oficjalny punkt końcowy \texttt{/v1/audio/transcriptions}, obsługujący pliki audio o rozmiarze do 25\,MB w formatach m.in. MP3, MP4, WAV, M4A i WebM \cite{openai_file_transcription}. Dostawca oferuje zarówno klasyczny model \texttt{whisper-1}, jak i nowsze generacje dedykowane do transkrypcji asynchronicznej (np. \texttt{gpt-transcribe}) oraz modele ze wsparciem dla automatycznej diarizacji mówców (\texttt{gpt-4o-transcribe-diarize}) \cite{openai_file_transcription}. Wariant chmurowy umożliwia również strumieniowanie zwracanych hipotez tekstowych w czasie przetwarzania gotowego pliku (\textit{file streaming}) \cite{openai_file_transcription}. \end{itemize} Model finansowy chmurowego API OpenAI opiera się na rozliczeniu per-skorzystanie za rzeczywisty czas nagrania audio (np. $0{,}0045$\,\$/min dla modelu \texttt{gpt-transcribe} oraz $0{,}006$\,\$/min dla \texttt{whisper-1}) bez stałych opłat abonamentowych \cite{openai_pricing}. Wśród funkcji wspomagających precyzję transkrypcji OpenAI API udostępnia opcję podawania kontekstu słownikowego poprzez parametr \texttt{prompt} (ograniczony w modelu Whisper do 224 tokenów) lub dedykowane parametry \texttt{keywords} oraz \texttt{languages}, które pozwalają na przekazywanie do modelu listy oczekiwanych słów kluczowych i skrótów branżowych \cite{openai_file_transcription}.

\subsection{Web Speech API (przeglądarkowe)}
\subsection{Porównanie dostępnych narzędzi:}
    \subsubsection{Dokładność transkrypcji}
    \subsubsection{Opóźnienia i czas przetwarzania}
    \subsubsection{Wymagania sprzętowe}
    \subsubsection{Koszty użytkowania}
    \subsubsection{Możliwość pracy offline}
    \subsubsection{Dostępność języka polskiego}

